\chapter*{Abstract}
\addcontentsline{toc}{chapter}{Abstract}

This work investigates the use of Abaqus-native error-guided mesh refinement for phase-field fracture simulations implemented with user elements. The point of departure is the staggered phase-field implementation of Moln\'ar and Gravouil and the mesh-refinement procedure proposed by Pandey and Kumar. The first objective is to reproduce reference fracture simulations without mesh refinement; the second is to implement the refinement procedure in Python while retaining the layered UEL/UMAT representation required by the available phase-field code; the third is to assess the additional requirements introduced by nonmatching state transfer and changing mesh topology.

A fixed-mesh Mode-I single-edge-crack calculation was reproduced before introducing adaptive refinement. The reconstructed response reached a maximum reaction force of $0.757778\,\mathrm{kN}$ at $u=0.005857\,\mathrm{mm}$. The project reference curve used for comparison gives approximately $0.758\,\mathrm{kN}$ at $u\approx0.005860\,\mathrm{mm}$, corresponding to differences of about $-0.029\%$ in peak load and $-0.051\%$ in peak displacement. This fixed-mesh calculation therefore provides the current quantitative baseline for the refinement study.

A Python-controlled two-job workflow was then implemented following Pandey and Kumar. A coarse pre-analysis supplies the Abaqus recovered-stress discretization indicator \texttt{MISESERI}; an Abaqus \texttt{RemeshingRule} and \texttt{adaptiveRemesh} operation generate a refined physical mesh; and the resulting connectivity is converted automatically into the coincident phase-field, mechanical and visualization layers used by the user-element formulation. A headless Abaqus qualification demonstrated genuine native remeshing and exposed an implementation detail that had initially overwritten the adaptively generated mesh by calling \texttt{generateMesh()} a second time.

The authoritative Task-5 production calculation (\texttt{1399632.mmaster02}) on the corrected sharp-slit $71{,}320$-element adaptive mesh completed technically to the requested displacement endpoint $u=0.010\,\mathrm{mm}$ in $35\,\mathrm{h}\,08\,\mathrm{min}$ of walltime. The calculation resolved the post-peak crack propagation branch to a final load drop of $94.4\%$ with a pure Mode-I horizontal crack path. However, the calculation failed the quantitative benchmark reproduction gates: the extracted peak load is $0.478203\,\mathrm{kN}$ at $u=0.004150\,\mathrm{mm}$ ($-36.89\%$ in peak force and $-29.14\%$ in peak displacement vs baseline), initial elastic stiffness is reduced by $-11.3\%$, and the native $1\%$ error target generated $71{,}320$ elements versus $\approx 13{,}941$ reported in the reference paper. A dedicated discrepancy audit confirmed that the UEL formulation, seam boundary conditions, and property vectors are mathematically sound, while the quantitative divergence reflects a combined structural compliance shift and earlier tip damage localization on the over-refined discretization. Task 5 remains open pending quantitative resolution.

In parallel, nonmatching transfer of displacement, phase field and history variables was investigated. Runtime state ingestion was implemented successfully, but topology-changing restart calculations showed reaction-force changes of approximately $26.3\%$ and $99.6\%$ during the relaxed restart stage for two tested crack-tip states. The current UEL also does not populate the Abaqus \texttt{ENERGY} array, so a directly comparable internal-energy continuity check is unavailable. Consequently, technically successful state installation is distinguished from mechanically consistent evolving-remesh continuation. The present shared-state implementation is qualified only for serial production execution because unsynchronised accesses to mutable \texttt{COMMON} data prevent a defensible threaded or MPI production claim.
